\documentclass[12pt]{article}

\usepackage{amsmath}
\usepackage{amssymb}
\usepackage{amsthm}
\usepackage[margin=1in]{geometry}
\usepackage{hyperref}
\usepackage{url}
\usepackage{booktabs}

\newtheorem{theorem}{Theorem}
\newtheorem{axiom}{Axiom}
\newtheorem{definition}{Definition}

\theoremstyle{remark}
\newtheorem*{scholium}{Scholium}

\title{The Inevitability of Tensorial Manifolds in Multi-Agent Ethics}
\author{Andrew Bond\\
\href{mailto:andrew.bond@sjsu.edu}{andrew.bond@sjsu.edu}}
\date{}

\begin{document}

\maketitle

\noindent\textbf{Keywords:} tensorial manifolds, multi-agent ethics, moral geometry, scalar insufficiency, tensor fields, metric tensor, differential ethics

\begin{abstract}
We provide a topological and algebraic proof that scalar-based ethical
frameworks (e.g., Utilitarianism) are insufficient for systems with non-trivial
symmetries. We demonstrate that for any moral judgment to be invariant under
``Reference Frame Transformations'' (shifts in cultural, individual, or temporal
perspectives), the underlying moral space must be a differentiable manifold
$\mathcal{M}$, and moral obligations must be represented as tensors of rank
$\geq 1$.
\end{abstract}

\bigskip
\hrule
\bigskip

\section*{I.\ Definitions and Axiomatization}

The tensorial framework presented here is developed fully in the Geometric Ethics manuscript~\cite{bond2026ge}, which provides the complete mathematical structure of moral reasoning on differentiable manifolds.

Let $\mathcal{M}$ be the \textbf{Moral Manifold}, a topological space
representing the set of all possible states of affairs.

\begin{enumerate}
    \item \textbf{Axiom of Perspective (Differentiability):} Every agent
    $\alpha$ possesses a local coordinate chart $\phi_\alpha : U_\alpha \to
    \mathbb{R}^n$, where $U_\alpha \subseteq \mathcal{M}$. This represents the
    agent's unique subjective view of the world.

    \item \textbf{Axiom of Consistency (Invariance):} A moral truth (e.g.,
    ``The Satisfaction of a Right'') must be a \textbf{geometric invariant}.
    That is, must be independent of the choice of coordinate chart
    $\phi_\alpha$.

    \item \textbf{The Obligation Vector:} An obligation $\mathbf{O}$ is defined
    as an element of the tangent space $T_p\mathcal{M}$, pointing toward a
    state of higher local value.

    \item \textbf{The Interest Covector:} An interest or ``right'' $\omega$ is
    a linear functional in the cotangent space $T_p^*\mathcal{M}$ that maps
    obligations to real-valued satisfaction.
\end{enumerate}

\bigskip
\hrule
\bigskip

\section*{II.\ The Proof of Scalar Insufficiency (The ``No-Go'' Theorem)}

\begin{theorem}
A rank-0 (scalar) ethical system cannot maintain consistency across
non-coincident agents.
\end{theorem}

\begin{proof}[Proof by Contradiction]
\begin{enumerate}
    \item Assume ethics is purely scalar, $V : \mathcal{M} \to \mathbb{R}$.

    \item Let Agent $\alpha$ and Agent $\beta$ observe the same act $A$. Due to
    their different positions in the manifold (different information, history, or
    values), they utilize transformation $\phi_\beta \circ \phi_\alpha^{-1}$.

    \item In a scalar system, if $\alpha$ measures $V(A) = v$, $\beta$ must
    also measure $V(A) = v$. However, if $A$ involves a direction (e.g.,
    ``returning a debt to $\beta$''), the value to $\alpha$ is not identical to
    the value to $\beta$.

    \item If $V_\alpha \neq V_\beta$, the scalar field is not invariant.
    Therefore, the system is \textbf{arbitrary} and fails the Axiom of
    Consistency.
\end{enumerate}
\end{proof}

\noindent\textbf{Conclusion:} Scalar ethics is only viable in a ``flat''
universe with a single agent, which is a trivial case.

\bigskip
\hrule
\bigskip

\section*{III.\ The Proof of Tensorial Necessity}

\begin{theorem}
The minimal mathematical structure satisfying Axioms 1 and 2 is a Tensor Field.
\end{theorem}

\begin{proof}
\begin{enumerate}
    \item To satisfy the Axiom of Consistency, we require an operation that
    yields a scalar which remains constant under the transformation
    $\phi_\beta \circ \phi_\alpha^{-1}$.

    \item According to the \textbf{Fundamental Theorem of Tensor Analysis}, the
    contraction of a tensor (vector) and a tensor (covector) is the unique
    linear operation that produces a coordinate-invariant scalar:
    \[
        S = \omega(\mathbf{O}) = \omega_i\, O^i
    \]

    \item By Definition 5.3, the transformation coefficients
    $\frac{\partial x^i}{\partial \bar{x}^j}$ and
    $\frac{\partial \bar{x}^j}{\partial x^i}$ cancel each other out during the
    contraction.

    \item Thus, $S = \omega(\mathbf{O})$ is the only ``Moral Fact'' that is
    objectively true for all agents simultaneously, even if they disagree on
    the ``direction'' of the obligation.
\end{enumerate}
\end{proof}

\bigskip
\hrule
\bigskip

\section*{IV.\ Optimization and the Metric Tensor}

To optimize moral outcomes, we must define the ``distance'' between the current
state $p$ and the ideal state $q$. This requires a \textbf{Metric Tensor}
$g_{ij}$.

\medskip

\noindent\textbf{The Proof of Bestness:} Any optimization of ``Moral Outcomes''
is equivalent to finding a geodesic on $(\mathcal{M}, g)$.

\begin{enumerate}
    \item The ``Ethical Stress'' of a system is represented by the
    \textbf{Ricci Curvature Tensor} $R_{ij}$.

    \item Optimization is defined as the minimization of the action:
    \[
        \mathcal{S} = \int \sqrt{g_{ij}\, dx^i\, dx^j}
    \]

    \item Because $R_{ij}$ accounts for the interaction between all dimensions
    of the manifold (the EM definitions and goals), any non-tensorial approach
    (like a simple list of priorities) will fail to see the ``curvature'' caused
    by competing interests, leading to systemic collapse (e.g., the ``tragedy
    of the commons'').
\end{enumerate}

\bigskip
\hrule
\bigskip

\section*{V.\ Final Scholium}

We have shown that:

\begin{enumerate}
    \item \textbf{Scalars} are blind to perspective (leading to
    totalitarianism).

    \item \textbf{Vectors} are blind to interaction (leading to chaos).

    \item \textbf{Tensors} are the unique mathematical structure that preserves
    \textbf{Relational Integrity}.
\end{enumerate}

Therefore, the \textbf{DEME (Distributed Ethical Moderation Engine)} is not
merely a software tool, but a computational implementation of the
\textbf{Universal Moral Metric}, now upgraded to DEME 3.0 with tensorial ethics support~\cite{bond2026erisml}. The ``Ethics Business'' is, in fact, the
engineering of the high-dimensional stability of the human species.

\begin{thebibliography}{9}

\bibitem{bond2026ge}
A.~H.~Bond, \emph{Geometric Ethics: The Mathematical Structure of Moral Reasoning}, v1.14, Department of Computer Engineering, San Jos\'e State University, Spring 2026.

\bibitem{bond2026erisml}
A.~H.~Bond, ``ErisML/DEME v3.0: A Library for Governed, Foundation-Model-Enabled Agents with Tensorial Ethics,'' \url{https://github.com/ahb-sjsu/erisml-lib}, 2026.

\end{thebibliography}

\end{document}
